%%%%%%%%%%%%%%%%%%%%%%%%%%%%%%%%%%%%%%%%%%%%%%%%%%%%%%%%%%%%%%%%%%%%%%
% Lớp 10 | Véctơ (chưa xét tọa độ) | Tích của một số với véctơ
% Dạng 0H5?3-1: Xác định 𝑘. và độ dài của nó
% Số câu: 13  (sinh tự động bởi tach_cau.py; từ nay thầy sửa trực tiếp file này)
%%%%%%%%%%%%%%%%%%%%%%%%%%%%%%%%%%%%%%%%%%%%%%%%%%%%%%%%%%%%%%%%%%%%%%

%%%%%%%%%%%%%%%% TRẮC NGHIỆM 4 PHƯƠNG ÁN %%%%%%%%%%%%%%%%

% nguồn: MaHoaNH_TracNghiem_Lop10
\begin{ex}%[0H5N3-1]%[ma:9ba2a79c]
	Cho tam  giác $ABC$. Gọi $M$ và $N$ lần lượt là trung điểm của $AB$ và $AC$. Khẳng định nào sau đây \textbf{sai}?
	 \choice
		 {$\vv{AB}=2\vv{AM}$}
		{$\vv{AC}=2\vv{NC}$} 
		{\True $\vv{BC}=-2\vv{MN}$} 
		 {$\vv{CN}=-\dfrac{1}{2}\vv{AC}$}
	\loigiai{ 
	\immini{Ta có $MN$ là đường trung bình của tam giác $ABC$  \\ $\Rightarrow\heva{&MN\parallel BC\\&BC=2MN.}$\\
		Mặt khác $\vv{MN}$  và $\vv{BC}$ là hai vectơ cùng chiều.\\
		 Do đó $\vv{BC}=2\vv{MN}$.}
		{ 
		\begin{tikzpicture}[scale=1,line join=round,font=\footnotesize,line cap=round,thick]
				\coordinate (A)  at (-0.5,3);
				\coordinate  (B) at (-2,0);
				\coordinate (C) at  (2.5,0);
				 \path ($(A)!0.5!(B)$) coordinate (M);
				 \path ($(A)!0.5!(C)$) coordinate (N);
				\draw (A)--(B)--(C)--cycle  (M)--(N);
				\foreach \i/\g in  {A/90,B/-90,C/-90,M/180,N/0}{\draw[fill=black](\i) circle (1pt)
					 ($(\i)+(\g:3mm)$) node[scale=1]{$\i$};}
			\end{tikzpicture}
		 }
	 }
\end{ex}

% nguồn: MaHoaNH_TracNghiem_Lop10
\begin{ex}%[0H5N3-1]%[ma:87855de2]
	Cho tam giác $ABC$. Gọi $I$ là trung điểm của $BC$.Khẳng định nào sau  đây đúng
	\choice 
		 {\True $\overrightarrow{BI}=\overrightarrow{IC}$}
		{$3\overrightarrow{BI}=2\overrightarrow{IC}$} 
		 {$\overrightarrow{BI}=\overrightarrow{2IC}$}
		{$\overrightarrow{2BI}=\overrightarrow{IC}$} 
	 \loigiai{
		\begin{center}
			\begin{tikzpicture}[=>stealth,line join=round,line cap=round,font=\footnotesize,scale=.7]
					 \path
						(0,0) coordinate  (A)
						 (1.5,3) coordinate (B)
						 (5,0) coordinate (C)
						($(B)!0.5!(C)$)  coordinate (I)
						 ($(A)!2/3!(I)$) coordinate (G);
					\draw  (A)--(B)--(C)--(A)
						 (A)--(I);
					\foreach \x/\g in{A/180,B/90,C/0,I/20,G/90} 
					\fill[black](\x)  circle (1pt)
						($(\x)+(\g:3mm)$)  node{$\x$};
				\end{tikzpicture}
		\end{center}
		Vì $I$  là trung điểm của $BC$ nên $BI=CI$ và $\overrightarrow{BI}$ cùng hướng với $\overrightarrow{IC}$ do đó hai vectơ $\overrightarrow{BI}$, $\overrightarrow{IC}$ bằng nhau hay $\overrightarrow{BI}=\overrightarrow{IC}$.
	} 
\end{ex}

% nguồn: MaHoaNH_TracNghiem_Lop10
\begin{ex}%[0H5H3-1]%[ma:74ff221a]
	Cho tam giác $ABC$. Gọi $I$ là trung điểm của $BC$.  Khẳng định nào sau đây đúng?
	 \choice
		 {\True $\overrightarrow{BI}=\overrightarrow{IC}$}
		{$3\overrightarrow{BI}=2\overrightarrow{IC}$} 
		{$\overrightarrow{BI}=\overrightarrow{2IC}$} 
		{$\overrightarrow{2BI}=\overrightarrow{IC}$} 
	\loigiai{ 
		Vì  $I$ là trung điểm của $BC$ nên $\overrightarrow{BI}=\overrightarrow{IC}$.
	 }
\end{ex}

% nguồn: MaHoaNH_TracNghiem_Lop10
\begin{ex}%[0H5H3-1]%[ma:5beb023d]
	Cho hai vectơ $\overrightarrow{a}$, $\overrightarrow{b}$ ngược hướng và $\left|\overrightarrow{a}\right|=5,\left|\overrightarrow{b}\right|=15$.  Với giá trị nào của $m$ thì $\overrightarrow{a}=m\overrightarrow{b}$?
	 \choice
		{$m=3$} 
		 {\True $m=-\dfrac{1}{3}$}
		 {$m=\dfrac{1}{3}$}
		{$m=-3$} 
	\loigiai{ 
		Do $\overrightarrow{a},\overrightarrow{b}$ ngược hướng nên  $m=-\dfrac{\left|\overrightarrow{a}\right|}{\left|\overrightarrow{b}\right|}=-\dfrac{5}{15}=-\dfrac{1}{3}$.
	} 
\end{ex}

% nguồn: MaHoaNH_TracNghiem_Lop10
\begin{ex}%[0H5H3-1]%[ma:cae5b2fe]
	Cho ba điểm phân biệt $A$, $B$, $C$. Nếu $\overrightarrow{AB}=-3\overrightarrow{AC}$ thì đẳng thức nào dưới đây  đúng?
	\choice 
		{$\overrightarrow{BC}=-4\overrightarrow{AC}$} 
		 {$\overrightarrow{BC}=-2\overrightarrow{AC}$}
		 {$\overrightarrow{BC}=2\overrightarrow{AC}$}
		 {\True $\overrightarrow{BC}=4\overrightarrow{AC}$}
	 \loigiai{
		 Ta có $\overrightarrow{AB}=-3\overrightarrow{AC} \Rightarrow \overrightarrow{BA} = 3\overrightarrow{AC}$ .\\
		 Do đó $\overrightarrow{BC} = \overrightarrow{BA} + \overrightarrow{AC} = 3\overrightarrow{AC}+\overrightarrow{AC} = 4\overrightarrow{AC}$.
	 }
\end{ex}

% nguồn: MaHoaNH_TracNghiem_Lop10
\begin{ex}%[0H5H3-1]%[ma:cb9e69d7]
	 Trên đường thẳng cho điểm $B$ nằm giữa hai điểm $A$ và $C$, với $AB=2a,AC=6a$. Đẳng thức nào sau đây đúng?
	 \choice
		{\True  $\overrightarrow{BC}=-2\overrightarrow{BA}$}
		 {$\overrightarrow{BC}=-2\overrightarrow{AB}$}
		{$\overrightarrow{BC}=4\overrightarrow{AB}$} 
		 {$\overrightarrow{BC}=\overrightarrow{AB}$}
	\loigiai{\immini{ 
		Có  $AB=2a\Rightarrow BC=4a=2AB$.
		\\Mà $\overrightarrow{BC}$; $\overrightarrow{AB}$ cùng hướng nên $$\overrightarrow{BC}=2\overrightarrow{AB}\Rightarrow \overrightarrow{BC}=-2\overrightarrow{BA}.$$}{\begin{tikzpicture}[scale=.7]
		\coordinate  [label=below:$A$](A) at (0,0);
		\coordinate [label=below:$B$](B)  at (2,0);
		 \coordinate [label=below:$C$](C) at (6,0);
		\foreach \diem  in {A,B,C} \fill (\diem) circle (1pt);
		 \draw[smooth](-2,0)--(8,0);
	\end{tikzpicture}
		 }
	} 
\end{ex}

% nguồn: MaHoaNH_TracNghiem_Lop10
\begin{ex}%[0H5Y3-1]%[ma:a117f4a6]
	Cho $\vec{a} = (2;-3), \vec{b} = (-1;4),  \vec{c} = \vec{a} - 2\vec{b}$. Khẳng định nào sau đây là \textbf{đúng}?
	\choice 
		{$\vec{c} =  (-4;-11)$}
		 {\True $\vec{c} = (4;-11)$}
		{$\vec{c} =  (4;11)$}
		 {$\vec{c} = (11;4)$}
	 \loigiai{
		Ta có $\vec{c} =  \vec{a} - 2\vec{b} = (2-2\cdot(-1);-3-2\cdot4) =(4; -11)$.
	} 
\end{ex}


%%%%%%%%%%%%%%%% ĐÚNG/SAI %%%%%%%%%%%%%%%%

% nguồn: MaHoaNH_DungSai_Lop10
\begin{ex}%[0H5N3-1]%[ma:47e728a6]
	Với hai véc-tơ $\vec{a}, \vec{b}$ và mọi số thực $h$ và $k$, xác định tính đúng sai của các mệnh  đề sau
	 \choiceTF
		 {\True $k(\vec{a}+\vec{b})=k\vec{a}+k\vec{b}$}
		 {\True $(h+k)\vec{a}=h\vec{a}+k\vec{a}$}
		{\True $h(k\vec{a})=(hk)\vec{a}$} 
		 {$(\vec{a}\cdot\vec{b})\cdot\vec{c}=\vec{a}\cdot(\vec{b}\cdot\vec{c})$}
	\loigiai{ 
		\begin{itemchoice}
			\itemch Theo tính chất, ta  có $k(\vec{a}+\vec{b})=k\vec{a}+k\vec{b}$.
			\itemch Theo tính chất,  ta có $(h+k)\vec{a}=h\vec{a}+k\vec{a}$.
			\itemch Theo  tính chất, ta có $h(k\vec{a})=(hk)\vec{a}$.
			\itemch Không có tính  chất kết hợp.
		\end{itemchoice}
	 }
\end{ex}

% nguồn: MaHoaNH_DungSai_Lop10
\begin{ex}%[0H5H3-1]%[ma:13d382cc]
	Cho tam giác $ABC$ và hai điểm $M, N$ thỏa mãn $\overrightarrow{BM}=\dfrac{2}{3}\overrightarrow{BC}$; $  \overrightarrow{CN}=-\dfrac{1}{2}\overrightarrow{AC}$. Xác định tính đúng sai của các mệnh đề sau
	\choiceTF 
		{$B, C, N$ thẳng  hàng}
		{\True $C, N, A$ thẳng  hàng}
		{\True $B, C, M$  thẳng hàng}
		{$C, B, A$  thẳng hàng}
	 \loigiai{
		\begin{itemchoice}
			\itemch $B, C, N$ không  thẳng hàng.
			 \itemch Do $ \overrightarrow{CN}=-\dfrac{1}{2}\overrightarrow{AC}$ nên $C, N, A$ thẳng hàng.
			\itemch Do $\overrightarrow{BM}=\dfrac{2}{3}\overrightarrow{BC}$ nên $B, C, M$  thẳng hàng.
			\itemch  $C,B,A$ không thẳng hàng.
	\end{itemchoice}}
\end{ex}


%%%%%%%%%%%%%%%% TRẢ LỜI NGẮN %%%%%%%%%%%%%%%%

% nguồn: MaHoaNH_TraLoiNgan_Lop10
\begin{bt}%[0H5H3-1]%[ma:57dc309f]
	Cho tam giác $ABC$. Trên cạnh $AB$ lấy điểm $M$ sao cho $2BA=5BM$. Gọi $G$ là trọng tâm tam giác $ABC$. Gọi $N$ là điểm trên $AC$ sao  cho $AN=xAC$. Tìm $x$, biết ba điểm $M$, $N$, $G$ thẳng hàng.
	 \shortans{0,75}
	\loigiai{\begin{center}
		\begin{tikzpicture}[scale=.7]
				\coordinate [label=left:$A$](A)  at (3,3);
				\coordinate [label=below:$B$](B)  at (0,0);
				\coordinate [label=below:$C$](C) at  (8,0);
				 \coordinate [label=left:$M$](M) at ($(A)! 3/5 ! (B)$);
				 \coordinate [label=right:$N$](N) at ($(A)! 3/4 ! (C)$);
				\coordinate  (P) at ($(B)!.5 ! (C)$);
				\coordinate [label=above right:$G$](G) at ($(A)! 2/3  ! (P)$);
				\foreach \diem in {A,B,C,M,N,P,G} \fill (\diem)  circle (1pt);
				 \draw[smooth](A)--(B)--(C)--(A)--(P)
					 (M)--(N);
			\end{tikzpicture}
	\end{center}
	Ta có 
	\begin{itemize}
		\item $\overrightarrow{MG}=\overrightarrow{AG}-\overrightarrow{AM}=\dfrac{1}{3}\left(\overrightarrow{AB}+\overrightarrow{AC}\right)-\dfrac{3}{5}\overrightarrow{AB}=\dfrac{1}{3}\overrightarrow{A C}-\dfrac{4}{15}\overrightarrow{AB}$. 
		\item $\overrightarrow{MN}=\overrightarrow{AN}-\overrightarrow{AM}=x \overrightarrow{AC}-\dfrac{3}{5}\overrightarrow{AB}$. 
	\end{itemize}
	Do $M$, $N$, $G$ thẳng hàng nên $\overrightarrow{MN}=k \overrightarrow{MG}\Rightarrow \dfrac{x}{\dfrac{1}{3}}=\dfrac{-\dfrac{3}{5}}{-\dfrac{4}{15}}\Rightarrow 3 x=\dfrac{9}{4}\Rightarrow x=\dfrac{3}{4}$. 
} 
\end{bt}

% nguồn: MaHoaNH_TraLoiNgan_Lop10
\begin{bt}%[0H5H3-1]%[ma:d4c81a41]
	Cho hình bình hành $ABCD$. Trên các đoạn thẳng $DC$, $AB$ theo thứ tự lấy  các điểm $M$, $N$ sao cho $DM=BN$. Gọi $P$ là giao điểm của $AM$, $DB$ và $Q$ là giao điểm của $CN$, $DB$. Nếu $\overrightarrow{DB}=k \overrightarrow{Q B}$ thì $k$ bằng bao nhiêu?
	\shortans{1} 
	 \loigiai{
		\begin{center}
			\begin{tikzpicture}[scale=.7]
					 \coordinate [label=left:$A$](A) at (3,3);
					\coordinate  [label=below:$B$](B) at (0,0);
					\coordinate [label=below:$C$](C)  at (8,0);
					\coordinate[label=right:$D$] (D) at  ($(C)-(B)+(A)$);
					\coordinate [label=right:$M$](M) at ($(D)!.7  ! (C)$);
					\coordinate [label=left:$N$](N) at  ($(B)!.7 ! (A)$);
					\draw[name  path=done] (A)--(M);
					\draw[name  path=dtwo,color=white] (B)--(D);
					\path [name intersections={of=done and  dtwo,by=P}];
					\draw (P) node[below]{$P$}; 
					\draw[name path=a]  (C)--(N);
					\draw[name path=b]  (B)--(D);
					\path [name intersections={of=a and b,by=Q}]; 
					\draw  (Q) node[below]{$Q$};
					\foreach \diem in {A,B,C,D,M,N,P,Q} \fill (\diem)  circle (1pt);
					 \draw[smooth](A)--(B)--(C)--(D)--(A);
				\end{tikzpicture}
		\end{center}
		Ta có $DM=BN \Rightarrow AN=MC$, mặt khác $AN$ song  song với $MC$ do đó tứ giác $ANCM$ là hình bình hành Suy ra $\overrightarrow{AM}=\overrightarrow{NC}$.\\
		Xét tam  giác $\triangle DMP$ và $\triangle BNQ$ ta có $DM=NB$ (giả thiết), $\widehat{PDM}=\widehat{QBN}$ (so le trong).\\
		Mặt khác $\widehat{D M P}=\widehat{A P B}$ (đối đỉnh) và $\widehat{A P Q}=\widehat{N Q B}$ (hai góc đồng vị) suy ra $\widehat{D M  P}=\widehat{B N Q}$.\\
		Do đó $\triangle DMP=\triangle BNQ$ (c.g.c) suy ra $DB=QB$.\\ 
		Dễ thấy $\overrightarrow{DB}$, $\overrightarrow{QB}$ cùng  hướng vì vậy $\overrightarrow{DB}=\overrightarrow{QB}$.
	} 
\end{bt}

% nguồn: MaHoaNH_TraLoiNgan_Lop10
\begin{bt}%[0H5V3-1]%[ma:5900d609]
	Cho tam giác $\triangle ABC$. Các điểm $I$, $J$ thỏa mãn hệ thức $\overrightarrow{AI}=\dfrac{1}{3} \overrightarrow{AB}$, $\overrightarrow{AI}=3\overrightarrow{AC}$. Biết $IC\parallel BJ$ thì $\overrightarrow{CI}=-\dfrac{a}{b} \overrightarrow{BJ}$, $\left(a, b \in \mathbb{N}^{*}\right)$ với $\dfrac{a}{b}$ là phân số tối  giản và $a, b \in \mathbb{Z}$. Tính $a+b$.
	\shortans{4} 
	\loigiai{ 
		Ta có  $\overrightarrow{AI}=\dfrac{1}{3} \overrightarrow{AB} \Leftrightarrow \overrightarrow{AC}+\overrightarrow{CI}=\dfrac{1}{3}\left(\overrightarrow{AC}+\overrightarrow{CB}\right) \Leftrightarrow \overrightarrow{CI}=-\dfrac{1}{3}\left(2\overrightarrow{AC}+\overrightarrow{BC}\right)$. \hfill(1)\\
		$\overrightarrow{A  J}=3 \overrightarrow{A C} \Leftrightarrow \overrightarrow{A B}+\overrightarrow{B J}=3 \overrightarrow{A B}+\overrightarrow{B C} \Leftrightarrow \overrightarrow{B J}=2 \overrightarrow{A B}+\overrightarrow{B C}=2 \overrightarrow{A C}+\overrightarrow{B C}$.\hfill(2)\\
		 Từ (1) và (2) $\Rightarrow \overrightarrow{CI}=-\dfrac{1}{3} \overrightarrow{BJ} \Rightarrow \dfrac{a}{b}=\dfrac{1}{3} \Rightarrow a+b=4$.
	} 
\end{bt}


%%%%%%%%%%%%%%%% TỰ LUẬN %%%%%%%%%%%%%%%%

% nguồn: MaHoaNH_TuLuan_Lop10
\begin{bt}%[0H5H3-1]%[ma:a8ef7375]
	Cho $\triangle A B C$  có điểm $D$, $I$ thỏa $3 \overrightarrow{DB}=2 \overrightarrow{D C}$, $\overrightarrow{I A}+3 \overrightarrow{I B}-2 \overrightarrow{I C}=\overrightarrow{0}$. Tính $k$ biết $\overrightarrow{A D}=k \overrightarrow{A I}$.
	\loigiai{\begin{center}
		\begin{center}
			\begin{tikzpicture}[scale=.7]
					\coordinate [label=above:$A$](A) at  (3,3);
					\coordinate  [label=below:$B$](B) at (0,0);
					\coordinate  [label=below:$C$](C) at (8,0);
					\coordinate [label=above right:$J$](J) at ($(A)! 2/5 ! (C)$); 
					\foreach  \diem in {A,B,C,J}\fill (\diem) circle (1.5pt);
					 \draw[smooth](A)--(B)--(C)--(A);
					 \draw[->](B)--(J);
				\end{tikzpicture}
		\end{center}
	\end{center}
	Ta có 
	\begin{eqnarray*}
		& & \overrightarrow{IA}+3 \overrightarrow{IB}-2  \overrightarrow{IC}=\overrightarrow{0}\\
		&\Leftrightarrow &  -\overrightarrow{AI}+3\left(\overrightarrow{AB}-\overrightarrow{AI}\right)-2\left(\overrightarrow{AC}-\overrightarrow{AI}\right)=\overrightarrow{0}\\
		&\Leftrightarrow  & -\overrightarrow{AI}+3 \overrightarrow{AB}-3 \overrightarrow{A I}-2 \overrightarrow{AC}+2 \overrightarrow{AI}=\overrightarrow{0}\\
		&\Leftrightarrow & 2 \overrightarrow{AI}=3\overrightarrow{AB}-2  \overrightarrow{AC}\\
		&\Leftrightarrow  & \overrightarrow{AI}=\dfrac{3}{2} \overrightarrow{AB}-\overrightarrow{AC}.
	\end{eqnarray*}
	Ta lại  có
	\begin{eqnarray*}
		& & 3\overrightarrow{DB}=2 \overrightarrow{DC}\\ 
		 &\Leftrightarrow & 3\left(\overrightarrow{AB}-\overrightarrow{A D}\right)=2\left(\overrightarrow{AC}-\overrightarrow{AD}\right)\\
		&\Leftrightarrow  & 3\overrightarrow{AB}-3\overrightarrow{AD}=2\overrightarrow{AC}-2\overrightarrow{AD}\\
		&\Leftrightarrow &\overrightarrow{AD}=3  \overrightarrow{AB}-2\overrightarrow{AC}.
	\end{eqnarray*}
	 Khi đó
	$\heva{&\overrightarrow{A I}=\dfrac{3}{2}\overrightarrow{AB}-\overrightarrow{AC} \\ &\overrightarrow{AD}=3\overrightarrow{AB}-2\overrightarrow{AC}} \Leftrightarrow \heva{&2 \overrightarrow{AI}=3\overrightarrow{AB}-2\overrightarrow{AC}  \\ &\overrightarrow{AD}=3\overrightarrow{AB}-2 \overrightarrow{AC}} \Leftrightarrow \overrightarrow{AD}=2 \overrightarrow{A I}$. \\Vậy $k=2$.
 }
\end{bt}
